% META: {"id":"1SPE-SUITES-EX-007","chapitre":"1SPE-SUITES","type_objet":"exercice","capacites":["1SPE-SUITES-C2","1SPE-SUITES-C5"],"capacites_codes":["C2","C5"],"methodes":["M2"],"parcours":1,"competences":["calculer"],"duree_min":10,"mode_creation":"ex_nihilo","sources_inspiration":[],"parametres_sympy":{"u0":100,"r":-3},"fichier_tex":"chapitres/1SPE-SUITES/exercices/1SPE-SUITES-EX-007.tex","corrige_tex":"chapitres/1SPE-SUITES/corriges/1SPE-SUITES-CO-007.tex","status":"approved","origin":"generated","status_history":["generated","audited","accepted","approved"],"release_acceptance":"RELEASE_OWNER_BATCH_ACCEPTANCE","acceptance_closure_digest":"sha256:9b3ccf9a81c5520fbb7e03b7b2d3e7bf2057a02834b6d908bf3be5f49f3b553f","acceptance_receipt_digest":"sha256:4fad86f0282e0fa4e0419cca1ad6379ae7af5a280c04a4a90880f90a1c044242"}
% BEGIN-VERIFY
% from sympy import *
% n = symbols('n', integer=True, nonnegative=True)
% u0 = 100
% r = -3
% u = u0 + n*r
% assert u.subs(n, 0) == 100
% assert u.subs(n, 1) == 97
% assert u.subs(n, 2) == 94
% assert u.subs(n, 20) == 40
% # u_p = u_0 + (p-0)*r => u_20 = 100 + 20*(-3) = 40
% assert 100 + 20*(-3) == 40
% END-VERIFY
\begin{exercice}{1SPE-SUITES-EX-007}{1}{10}
Sophie étudie une suite arithmétique $(u_n)$ de premier terme $u_0 = 100$ et de raison $r = -3$.
\begin{enumerate}
  \item Écrire la formule du terme général $u_n$ en fonction de $n$.
  \item Calculer $u_1$ et $u_2$.
  \item Calculer $u_{20}$.
  \item La suite est-elle croissante ou décroissante ? Justifier à partir de la raison.
\end{enumerate}
\end{exercice}
